\chapter{Metodología de desarrollo}

En este capítulo se describe la metodología seguida durante el desarrollo y validación de las distintas implementaciones de FlashAttention. En primer lugar, se presenta el núcleo algorítmico por bloques de FlashAttention y se distingue de las mejoras de particionado introducidas por FlashAttention-2. Posteriormente, se detalla la implementación desarrollada directamente en CUDA, incluyendo las principales decisiones de diseño y optimizaciones aplicadas.

A continuación, se describe la implementación realizada con Triton y el proceso de \textit{autotuning} utilizado para seleccionar sus parámetros de ejecución. Finalmente, se presenta la metodología empleada para verificar la correctitud de los resultados y para detectar y depurar errores durante el desarrollo de los kernels.

Los kernels implementados en CUDA y Triton, así como la implementación de referencia basada en SDPA, se evalúan bajo una configuración común: dimensión de cabeza igual a 128, longitud de secuencia múltiplo de 16, un único elemento por batch y una única cabeza de atención. En todos los casos se emplea atención no causal.


\section{FlashAttention y FlashAttention-2}

El núcleo algorítmico utilizado como punto de partida es FlashAttention: atención exacta por bloques con actualización incremental del \textit{softmax}, de forma que no sea necesario materializar completamente la matriz \(QK^T\) en memoria global \cite{dao2022flashattention}. El \textit{softmax online} empleado para mantener el máximo y el término de normalización durante el recorrido incremental se apoya en la formulación de Milakov y Gimelshein \cite{milakov2018onlineSoftmax}.

FlashAttention-2 conserva este principio, pero añade cambios relevantes de particionado y paralelización. En particular, reduce trabajo no matricial, paraleliza la atención de una misma cabeza entre varios bloques y modifica la distribución del trabajo entre \textit{warps} para disminuir comunicación y aumentar la utilización de la GPU \cite{dao2023flashattention2}. La implementación CUDA de este TFG \textbf{no pretende reproducir íntegramente el particionado de FlashAttention-2}; utiliza su algoritmo por bloques y varias de sus ideas como referencia, pero adopta un \textit{mapping} propio que posteriormente se analiza de forma experimental.

Por esta razón, el siguiente pseudocódigo debe interpretarse como el \textbf{núcleo algorítmico común de FlashAttention}, no como una reproducción completa de todas las decisiones de paralelización de FlashAttention-2.

\begin{algorithm}[H]
\caption{Núcleo algorítmico por bloques de FlashAttention}
\label{alg:flashattention2}
\begin{algorithmic}[1]

\Require \(Q,K,V \in \mathbb{R}^{N\times d}\)
\Require \(B_r\): número de filas por bloque de \(Q\)
\Require \(B_c\): número de filas por bloque de \(K\) y \(V\)
\Ensure \(O = \operatorname{softmax}(QK^T/\sqrt{d})V\)

\State \(T_r \gets \lceil N/B_r \rceil\)
\State \(T_c \gets \lceil N/B_c \rceil\)

\State Dividir \(Q\) en bloques \(Q_1,\ldots,Q_{T_r}\), con \(Q_i \in \mathbb{R}^{B_r\times d}\)
\State Dividir \(K\) y \(V\) en bloques \(K_1,\ldots,K_{T_c}\) y \(V_1,\ldots,V_{T_c}\), con \(K_j,V_j \in \mathbb{R}^{B_c\times d}\)

\For{\(i \gets 1,\ldots,T_r\)}
\State \(O_i \gets 0 \in \mathbb{R}^{B_r\times d}\)
\State \(m_i \gets -\infty \in \mathbb{R}^{B_r}\)
\State \(l_i \gets 0 \in \mathbb{R}^{B_r}\)
\For{\(j \gets 1,\ldots,T_c\)}
    \State \(S_{ij} \gets Q_iK_j^T/\sqrt{d}\)
    \State \(m_i^{new} \gets \max(m_i,\operatorname{rowmax}(S_{ij}))\)
    \State \(P_{ij} \gets \exp(S_{ij}-m_i^{new})\)
    \State \(\alpha_i \gets \exp(m_i-m_i^{new})\)
    \State \(l_i^{new} \gets \alpha_i \odot l_i + \operatorname{rowsum}(P_{ij})\)
    \State \(O_i \gets \alpha_i \odot O_i + P_{ij}V_j\)
    \State \(m_i \gets m_i^{new}\)
    \State \(l_i \gets l_i^{new}\)
\EndFor
\State \(O_i \gets O_i \oslash l_i\)
\State Escribir \(O_i\) en \(O\)
\EndFor
\end{algorithmic}
\end{algorithm}

En el pseudocódigo anterior, \(\operatorname{rowmax}\) y \(\operatorname{rowsum}\) calculan respectivamente el máximo y la suma de cada fila. Los vectores \(m_i\) y \(l_i\) almacenan el estado necesario para la normalización incremental. El pseudocódigo expresa dependencias matemáticas; no fija por sí mismo el nivel de la jerarquía de memoria ni el reparto concreto entre bloques y \textit{warps}.

\subsection{Algoritmo CUDA}

La implementación CUDA distribuye el cálculo de la atención a nivel de \textit{warp}. Cada \textit{warp} es responsable de un bloque de 16 filas consecutivas de la matriz de consultas \(Q\), de forma que un bloque CUDA contiene varios \textit{warps} que procesan de manera independiente distintos bloques de consultas. La posición global del \textit{warp} determina las filas de \(Q\) y de la salida \(O\) que le corresponden, mientras que su identificador dentro del bloque se emplea para cooperar en la transferencia de \(K\) y \(V\) a memoria compartida. De esta manera, cada \textit{warp} mantiene en registros su bloque de consultas y el acumulador correspondiente a 16 filas de la salida.

Las matrices \(K\) y \(V\) se recorren en bloques de \(16\times128\) elementos. Para ocultar parcialmente la latencia de memoria global se utiliza un \textit{pipeline} con múltiples buffers en memoria compartida. Mientras un bloque de \(K\) y \(V\) está siendo consumido por los \textit{warps}, bloques posteriores se transfieren desde memoria global mediante instrucciones asíncronas \texttt{cp.async}. Los buffers se organizan de forma circular, de manera que el buffer utilizado para el bloque \(i\) puede reutilizarse para almacenar el bloque \(i+\texttt{buffer\_num}\) una vez que todos los \textit{warps} han terminado de consumirlo. La carga de cada bloque de \(K\) y \(V\) se realiza de forma cooperativa entre todos los \textit{warps} del bloque CUDA. Como consecuencia, todos ellos avanzan de forma prácticamente \textit{lockstep} a través del bucle principal: antes de consumir un nuevo bloque deben esperar a que las transferencias asíncronas correspondientes hayan finalizado y sincronizarse mediante una barrera a nivel de bloque. De forma análoga, antes de reutilizar un buffer de memoria compartida, todos los \textit{warps} deben haber terminado de consumir su contenido. Esta sincronización garantiza la corrección del acceso cooperativo a memoria compartida, aunque reduce el grado de independencia temporal entre los distintos \textit{warps} del bloque.
La transferencia desde memoria compartida al fichero de registros constituye un segundo nivel de \textit{buffering}; las instrucciones empleadas se corresponden con las primitivas documentadas en la ISA PTX de NVIDIA \cite{nvidiaPTX}. Los fragmentos de \(K\) y \(V\) se cargan mediante \texttt{ldmatrix.sync.aligned.m8n8.x4.shared.b16}, que permite obtener simultáneamente cuatro matrices \(8\times8\) desde memoria compartida y distribuirlas directamente entre los registros del \textit{warp}. El código mantiene un pequeño \textit{pipeline} entre estas cargas y las operaciones matriciales: mientras un par de fragmentos se utiliza en los Tensor Cores, se carga desde memoria compartida el siguiente par que será consumido.

Los productos matriciales se realizan mediante instrucciones \texttt{mma.sync} de tipo \(m16n8k16\), con operandos BF16 y acumulación FP32. Cada instrucción calcula un producto \(16\times16\) por \(16\times8\), produciendo un fragmento \(16\times8\) en precisión simple. Para calcular \(QK^\mathsf{T}\), un bloque lógico de puntuaciones de tamaño \(16\times16\) se representa mediante dos fragmentos \(16\times8\). Tras aplicar el escalado y el \textit{softmax} online, las probabilidades se convierten nuevamente a BF16 y se utilizan como operando izquierdo en los productos \(PV\). El acumulador de salida permanece en FP32 durante todo el recorrido de \(K\) y \(V\), realizándose la conversión a BF16 únicamente al escribir el resultado final.
Para mejorar el acceso posterior mediante \texttt{ldmatrix}, los bloques de \(K\) y \(V\) se almacenan en memoria compartida utilizando un patrón \textit{swizzled} con granularidad de ocho elementos BF16 (16 bytes). La transformación aplica un XOR entre la posición lógica del segmento y los bits bajos de la fila, redistribuyendo los accesos entre los bancos de memoria compartida. Al cargar posteriormente los fragmentos mediante \texttt{ldmatrix}, se aplica la misma transformación para reconstruir la dirección física correspondiente. De este modo, la disposición en memoria global puede mantenerse contigua y coalescente mientras la disposición en memoria compartida se adapta al patrón de acceso de las cargas matriciales.

En conjunto, el kernel emplea dos niveles explícitos de movimiento de datos: un \textit{pipeline} memoria global--memoria compartida mediante \texttt{cp.async}, y un segundo \textit{pipeline} memoria compartida--registros mediante \texttt{ldmatrix}. Sobre los fragmentos residentes en registros se ejecutan las instrucciones \texttt{mma.sync} de los Tensor Cores, manteniendo tanto los estados del \textit{softmax} online como el acumulador de salida en registros durante toda la operación.

\begin{algorithm}[H]
\caption{FlashAttention CUDA}
\footnotesize
\begin{algorithmic}[1]

\State Asignar a cada \textit{warp} 16 filas consecutivas de \(Q\)
\State Cargar el bloque de \(Q\) correspondiente en registros
\State Inicializar acumulador \(O \gets 0\), máximo de fila \(m \gets -\infty\) y denominador \(l \gets 0\)

\For{\(b \gets 0\) hasta \(\texttt{buffer\_num}-1\)}
\State Cargar cooperativamente \(K_b\) y \(V_b\) de memoria global a memoria compartida mediante \texttt{cp.async}
\State Aplicar \textit{swizzle} XOR durante la escritura en memoria compartida
\EndFor

\For{cada bloque \(i\) de 16 filas de \(K\) y \(V\)}

\State Esperar al buffer correspondiente mediante \texttt{cp.async.wait\_group}
\State Sincronizar todos los \textit{warps} del bloque
\Comment{Los \textit{warps} avanzan prácticamente en \textit{lockstep}}

\State \(S \gets 0\)

\State Cargar primeros fragmentos de \(K_i\) desde memoria compartida a registros con \texttt{ldmatrix}

\For{cada fragmento de dimensión \(16\times16\) en la dimensión de cabeza}
    \State Cargar anticipadamente el siguiente fragmento de \(K_i\) mediante \texttt{ldmatrix}
    \State \(S \gets S + QK_i^\mathsf{T}\) mediante \texttt{mma.sync.m16n8k16}
\EndFor

\State \(S \gets S/\sqrt{d}\)

\State \(m_{\mathrm{new}} \gets \max(m,\operatorname{rowmax}(S))\)
\State \(\alpha \gets \exp(m-m_{\mathrm{new}})\)
\State \(P \gets \exp(S-m_{\mathrm{new}})\)
\State \(l_{\mathrm{new}} \gets \alpha l + \operatorname{rowsum}(P)\)

\State Convertir \(P\) de FP32 a BF16 para utilizarlo como operando de los Tensor Cores

\State \(O \gets \alpha O\)

\State Cargar primeros fragmentos de \(V_i\) desde memoria compartida a registros mediante \texttt{ldmatrix}

\For{cada fragmento \(16\times8\) de \(V_i\)}
    \State Cargar anticipadamente el siguiente fragmento de \(V_i\)
    \State \(O \gets O + PV_i\) mediante \texttt{mma.sync.m16n8k16}
\EndFor

\State \(m \gets m_{\mathrm{new}}\)
\State \(l \gets l_{\mathrm{new}}\)

\State Sincronizar todos los \textit{warps} antes de reutilizar el buffer

\If{existen bloques futuros de \(K\) y \(V\)}
    \State Reutilizar el buffer consumido
    \State Cargar asíncronamente el siguiente bloque \(K_{i+\texttt{buffer\_num}},V_{i+\texttt{buffer\_num}}\)
\EndIf

\EndFor

\State \(O \gets O/l\)

\State Redistribuir los elementos de \(O\) entre los hilos del \textit{warp}
\State Escribir \(O\) en memoria global mediante stores coalescentes

\end{algorithmic}
\end{algorithm}

\subsection{FlashAttention: Triton}

La implementación en Triton sigue la misma descomposición algorítmica que las versiones anteriores. Cada instancia del programa Triton procesa un bloque de \(B_r\) filas consecutivas de la matriz de consultas \(Q\) y recorre secuencialmente la dimensión de secuencia de \(K\) y \(V\) en bloques de \(B_c\) filas. De esta forma, cada programa mantiene localmente el acumulador de salida y los parámetros necesarios para calcular el \textit{softmax online}, evitando materializar la matriz completa \(QK^\mathsf{T}\) en memoria global.

Dado que la dimensión de cabeza se fija en \(d=128\), cada bloque de consultas tiene dimensiones \(B_r\times128\), mientras que los bloques de claves y valores tienen dimensiones \(B_c\times128\). Para cada bloque de \(K\) se calcula una matriz temporal de puntuaciones de dimensiones \(B_r\times B_c\). Esta matriz únicamente permanece durante la iteración actual del algoritmo y se utiliza inmediatamente para actualizar el \textit{softmax} y el acumulador de salida.

El producto matricial se expresa mediante la operación \texttt{tl.dot}. Con operandos BF16 sobre arquitecturas que disponen de Tensor Cores, el compilador de Triton genera las instrucciones matriciales correspondientes y mantiene la acumulación en FP32.

La implementación se parametriza mediante cuatro parámetros principales:

\begin{itemize}
\item \(B_r\): número de filas de \(Q\) procesadas por cada instancia del programa.
\item \(B_c\): número de filas de \(K\) y \(V\) procesadas en cada iteración.
\item \texttt{num\_warps}: número de \textit{warps} utilizados para ejecutar una instancia del programa.
\item \texttt{num\_stages}: número de etapas utilizadas por el \textit{pipeline} de cargas y cómputo generado por Triton.
\end{itemize}

La dimensión de cabeza no forma parte del espacio de búsqueda, ya que se mantiene fija en 128 para todos los experimentos. De forma análoga, el tamaño de batch y el número de cabezas se fijan ambos a uno. Por tanto, el proceso de \textit{autotuning} explora únicamente combinaciones de los cuatro parámetros anteriores.

El valor de \(B_r\) determina la cantidad de filas de salida calculadas simultáneamente y, por tanto, afecta directamente al tamaño del acumulador mantenido por cada programa. El valor de \(B_c\) determina el tamaño de los bloques temporales de puntuaciones y la granularidad con la que se recorren \(K\) y \(V\). Valores elevados de ambos parámetros pueden incrementar la reutilización de datos y el trabajo realizado por programa, pero también aumentan el consumo de registros y memoria compartida. Por su parte, \texttt{num\_warps} controla el paralelismo dentro de cada programa, mientras que \texttt{num\_stages} modifica la profundidad del \textit{pipeline} utilizado para solapar transferencias de memoria y cómputo.

Las distintas combinaciones se evalúan mediante el mecanismo de \textit{autotuning} de Triton, seleccionándose para cada configuración aquella que presenta el menor tiempo de ejecución. De este modo se exploran automáticamente distintas relaciones entre tamaño de bloque, paralelismo y profundidad del \textit{pipeline}.

El funcionamiento del kernel puede resumirse mediante el Algoritmo~\ref{alg:flashattention_triton}.

\begin{algorithm}[H]
\caption{FlashAttention Triton}
\label{alg:flashattention_triton}
\begin{algorithmic}[1]

\Require \(Q,K,V\in\mathbb{R}^{N\times128}\)
\Require \(B_r\): filas de \(Q\) procesadas por programa
\Require \(B_c\): filas de \(K\) y \(V\) procesadas por iteración
\Ensure \(O=\operatorname{softmax}(QK^\mathsf{T}/\sqrt{128})V\)

\State Asignar a cada programa un bloque \(Q_i\in\mathbb{R}^{B_r\times128}\)
\State Cargar \(Q_i\)
\State \(O_i \gets 0 \in \mathbb{R}^{B_r\times128}\) en FP32
\State \(m_i \gets -\infty \in \mathbb{R}^{B_r}\)
\State \(l_i \gets 0 \in \mathbb{R}^{B_r}\)

\For{cada bloque \(K_j,V_j\in\mathbb{R}^{B_c\times128}\)}

\State Cargar \(K_j\)
\State \(S_{ij} \gets \texttt{tl.dot}(Q_i,K_j^\mathsf{T})/\sqrt{128}\)

\State \(m_i^{new}
    \gets
    \max\left(
        m_i,
        \operatorname{rowmax}(S_{ij})
    \right)\)

\State \(\alpha_i
    \gets
    \exp(m_i-m_i^{new})\)

\State \(P_{ij}
    \gets
    \exp(S_{ij}-m_i^{new})\)

\State \(l_i^{new}
    \gets
    \alpha_i\odot l_i
    +
    \operatorname{rowsum}(P_{ij})\)

\State \(O_i \gets \alpha_i\odot O_i\)

\State Cargar \(V_j\)
\State \(O_i
    \gets
    O_i
    +
    \texttt{tl.dot}(P_{ij},V_j)\)

\State \(m_i \gets m_i^{new}\)
\State \(l_i \gets l_i^{new}\)

\EndFor

\State \(O_i \gets O_i\oslash l_i\)
\State Convertir \(O_i\) al formato de salida
\State Escribir \(O_i\) en memoria global

\end{algorithmic}
\end{algorithm}

El pseudocódigo describe el flujo algorítmico del programa Triton; la planificación concreta de las operaciones y del movimiento de datos depende del código generado para la configuración seleccionada.

\chapter{Entorno de prueba}

En este capítulo se describe el entorno hardware y software utilizado para la evaluación experimental de las distintas implementaciones de FlashAttention. Todas las mediciones presentadas en los capítulos posteriores se han realizado sobre la misma plataforma, con el objetivo de que las comparaciones entre CUDA, Triton y las implementaciones proporcionadas por PyTorch se efectúen bajo condiciones equivalentes.

El entorno de pruebas corresponde a una máquina virtual proporcionada a través de la plataforma Vast.ai. Esta plataforma permite acceder de forma remota a sistemas equipados con GPUs de altas prestaciones. En concreto, para los experimentos se seleccionó una instancia equipada con una GPU NVIDIA A100-SXM4 de 40 GB. Al tratarse de un entorno virtualizado, algunos parámetros de bajo nivel del sistema anfitrión, como el estado de \textit{boost} del procesador, no son accesibles desde la máquina virtual.

\section{Hardware}

Los experimentos se han ejecutado sobre una GPU NVIDIA A100-SXM4 con 40 GB de memoria. Esta GPU pertenece a la arquitectura Ampere y presenta capacidad de cómputo 8.0 (\texttt{sm\_80}). El dispositivo dispone de 108 \textit{Streaming Multiprocessors} (SM) y está específicamente orientado a cargas de trabajo de cómputo de alto rendimiento.

La Tabla~\ref{tab:gpu_entorno} resume las principales características del dispositivo empleado.

\begin{table}[H]
\centering
\caption{Características de la GPU utilizada durante los experimentos.}
\label{tab:gpu_entorno}
\begin{tabular}{ll}
\hline
\textbf{Característica} & \textbf{Valor} \\
\hline
Proveedor de la instancia & Vast.ai \\
GPU & NVIDIA A100-SXM4-40GB \\
Arquitectura & Ampere \\
Capacidad de cómputo & 8.0 (\texttt{sm\_80}) \\
Número de SM & 108 \\
Memoria & 40 GB \\
Límite de potencia & 400 W \\
Frecuencia máxima de los SM & 1410 MHz \\
Frecuencia de memoria & 1215 MHz \\
\hline
\end{tabular}
\end{table}

La máquina virtual dispone además de recursos de CPU pertenecientes a un procesador AMD EPYC 7343. La instancia utilizada expone 7 CPU lógicas al sistema invitado, con una frecuencia observada de aproximadamente 3.19 GHz. La máquina dispone de aproximadamente 49 GiB de memoria principal. Debido a la virtualización, la información relativa al estado de \textit{boost} del procesador no es expuesta por el sistema y, por tanto, no se controla explícitamente durante los experimentos.

\section{Entorno software}

El sistema operativo utilizado es Ubuntu 22.04.5 LTS, con kernel Linux 6.8.0-138-generic y arquitectura x86-64.

Las implementaciones CUDA se han compilado utilizando CUDA Toolkit 13.0 y \texttt{nvcc} 13.0.88. El código C++ auxiliar se compila mediante GCC 11.4.0 y la generación del proyecto se realiza con CMake 4.4.3.

La Tabla~\ref{tab:software_entorno} recoge las principales versiones utilizadas.

\begin{table}[H]
\centering
\caption{Entorno software utilizado durante los experimentos.}
\label{tab:software_entorno}
\begin{tabular}{ll}
\hline
\textbf{Componente} & \textbf{Versión} \\
\hline
Ubuntu & 22.04.5 LTS \\
Kernel Linux & 6.8.0-138-generic \\
Driver NVIDIA & 580.173.02 \\
CUDA Toolkit & 13.0 \\
\texttt{nvcc} & 13.0.88 \\
GCC / G++ & 11.4.0 \\
CMake & 4.4.3 \\
Python & 3.14.7 \\
PyTorch & 2.13.0+cu130 \\
CUDA de PyTorch & 13.0 \\
cuDNN & 9.2.0 \\
Triton & 3.7.1 \\
Nsight Compute & 2025.3.1 \\
\hline
\end{tabular}
\end{table}

PyTorch detecta correctamente la GPU como una A100 con capacidad de cómputo 8.0. Tanto la implementación Triton como las implementaciones de referencia basadas en SDPA se ejecutan dentro del mismo entorno de Python, evitando diferencias derivadas de utilizar distintas versiones del \textit{runtime} de CUDA.

\section{Metodología de medida}

Las distintas implementaciones se evalúan utilizando las mismas dimensiones de entrada, el mismo tipo de datos y la misma política de generación de entradas. Se considera un único elemento por \textit{batch}, una única cabeza de atención, dimensión de cabeza \(d=128\) y atención no causal. Los tensores \(Q\), \(K\) y \(V\) utilizan formato BF16.

Todos los tests de rendimiento utilizan la misma semilla pseudoaleatoria, fijada a 0 mediante \path{torch.manual_seed(0)} y \path{torch.cuda.manual_seed_all(0)} antes de generar las entradas. De esta forma, CUDA, Triton, SDPA y cuDNN parten de entradas reproducibles construidas bajo el mismo criterio para una misma longitud de secuencia. Esta decisión evita introducir una fuente de variación innecesaria al comparar implementaciones.

Para las medidas temporales se realiza inicialmente una ejecución de calentamiento. Esta primera ejecución permite excluir operaciones que ocurren únicamente al comienzo de la prueba, como la compilación JIT de los kernels Triton o la inicialización de determinados componentes del entorno CUDA.

Posteriormente, cada configuración se ejecuta 100 veces. El tiempo de ejecución se obtiene mediante eventos CUDA, evitando incluir en la medida el tiempo empleado por Python o por otras operaciones realizadas en CPU. Para cada longitud de secuencia se calcula la media y la desviación típica de los tiempos obtenidos. Las cuatro rutas se evalúan con el mismo conjunto de longitudes de secuencia y el mismo número de repeticiones.

Las dos rutas de PyTorch se fuerzan de forma explícita mediante el contexto \path{torch.nn.attention.sdpa_kernel}: una ejecución selecciona \path{SDPBackend.FLASH_ATTENTION} y la otra \path{SDPBackend.CUDNN_ATTENTION}. De este modo, el benchmark no depende de la selección automática del backend realizada por PyTorch \cite{pytorchSDPA}. En Triton, la primera llamada asociada a una configuración se realiza fuera de la región temporizada; así, la compilación JIT y las evaluaciones necesarias para el \textit{autotuning} no forman parte de las 100 repeticiones medidas \cite{tritonAutotune}.

La instancia de Vast.ai no fija manualmente las frecuencias de GPU. Por tanto, las medidas pueden estar sometidas a la política dinámica de frecuencia y a cierta variabilidad propia de un entorno virtualizado. Para reducir su impacto, todas las implementaciones se miden en la misma instancia y sesión experimental y se reportan media y desviación típica. Las desviaciones observadas son pequeñas respecto a las diferencias entre implementaciones, pero esta condición se conserva como limitación de validez externa.

Para el análisis detallado de los kernels se utiliza NVIDIA Nsight Compute 2025.3.1 \cite{nsightCompute}. Los tests de perfilado realizan un calentamiento previo y delimitan mediante NVTX una única ejecución del kernel de interés. La instrumentación de Nsight Compute puede introducir \textit{replay} y alterar la duración observada, por lo que los tiempos del perfilador no se utilizan como benchmark principal. Las métricas recogidas se emplean únicamente para explicar la ocupación, el tráfico de memoria, el uso de registros y memoria compartida, la actividad de Tensor Cores y las causas de espera de los \textit{warps}.

\chapter{Tests}

Durante el desarrollo se utilizó un test funcional para comprobar la corrección numérica y un conjunto separado de tests no funcionales para estudiar el rendimiento mediante medidas temporales y herramientas de \textit{profiling}.

Los tests se implementan utilizando \texttt{pytest}. Las pruebas correspondientes a CUDA, Triton y las implementaciones de referencia de PyTorch se mantienen separadas, permitiendo ejecutar de forma independiente las diferentes familias de tests.

\section{Test funcional}

La validación funcional se concentra en una única prueba de referencia. Para cada implementación se generan tensores pseudoaleatorios \(Q\), \(K\) y \(V\) en BF16 y se compara la salida obtenida con la producida por \path{scaled_dot_product_attention} de PyTorch bajo la misma configuración.

Para garantizar la reproducibilidad se fija la semilla a 0 antes de generar las entradas. La comparación se realiza mediante \texttt{torch.allclose}, utilizando tolerancias relativa y absoluta de \(2\cdot10^{-2}\). No se exige igualdad bit a bit porque las implementaciones pueden ejecutar las operaciones en distinto orden y utilizar BF16 en puntos diferentes del cálculo.

Este test se ejecuta después de los cambios relevantes del kernel y actúa como comprobación automática de que las optimizaciones no alteran el resultado numérico esperado. Se utiliza una única semilla y una tolerancia fija; por tanto, la validación demuestra compatibilidad numérica para los casos ejercitados, pero no constituye una caracterización exhaustiva del error numérico. El objetivo del trabajo no es estudiar estrategias de validación de atención, por lo que no se incorporan campañas multi-semilla ni métricas específicas de error máximo o medio.

\section{Tests no funcionales}

Los tests no funcionales no verifican directamente el valor de la salida, sino que se utilizan para caracterizar el rendimiento de las diferentes implementaciones. Se distinguen dos tipos: tests de tiempo y tests destinados al \textit{profiling}.

\subsection{Tests de tiempo}

Los tests de tiempo permiten estudiar cómo varía el coste de la operación al aumentar la longitud de secuencia. Se utilizan las mismas dimensiones y tipos de datos para CUDA, Triton y SDPA, de forma que los resultados puedan compararse directamente.

Todos los tests temporales inicializan los generadores pseudoaleatorios con la misma semilla, 0, antes de generar las entradas. Este criterio se mantiene para CUDA, Triton, SDPA y cuDNN.

Para cada longitud de secuencia se generan los tensores \(Q\), \(K\) y \(V\) en BF16. Antes de comenzar la medida se realiza al menos una ejecución de calentamiento. Esta ejecución es especialmente importante para Triton, ya que permite excluir de la medida el tiempo correspondiente a la compilación JIT del kernel.

El tiempo se mide utilizando eventos CUDA. Para cada repetición se registra un evento inmediatamente antes del lanzamiento de la implementación y otro inmediatamente después. De esta forma, la medida corresponde al trabajo realizado en la GPU y no incluye el tiempo empleado por Python en preparar la llamada.

Cada configuración se ejecuta múltiples veces. Si los tiempos obtenidos son

$$
t_1,t_2,\ldots,t_R,
$$

el tiempo utilizado como resultado es la media

$$
\bar{t}
=
\frac{1}{R}
\sum_{i=1}^{R}t_i,
$$

acompañada de su desviación típica

$$
\sigma =
\sqrt{
\frac{1}{R-1}
\sum_{i=1}^{R}
(t_i-\bar{t})^2
}.
$$

Por tanto, para cada implementación se obtiene un conjunto de valores de la forma

$$
(N,\bar{t},\sigma),
$$

que permite estudiar tanto el rendimiento medio como la variabilidad de las mediciones al incrementar la longitud de secuencia.

Los tests destinados a estas medidas se identifican mediante un marcador específico de \texttt{pytest}, lo que permite ejecutarlos independientemente del test funcional. De esta manera, las pruebas de rendimiento pueden lanzarse mediante una única batería común para CUDA, Triton y SDPA.

\subsection{Tests de profiling}

Para estudiar las causas de las diferencias de rendimiento se utilizan tests específicos de \textit{profiling}. A diferencia de los tests temporales, estos tests realizan únicamente las ejecuciones necesarias para que una herramienta externa pueda recoger métricas del kernel.

Antes de la ejecución analizada se realiza una llamada de calentamiento seguida de una sincronización de CUDA. Posteriormente se lanza una única ejecución del kernel de interés. La región correspondiente se delimita mediante anotaciones NVTX, lo que permite identificarla desde herramientas de análisis externas y evitar la captura de las operaciones utilizadas para preparar los tensores o realizar el calentamiento.

Los perfiles se obtienen utilizando NVIDIA Nsight Compute. Entre las métricas estudiadas se encuentran la ocupación de los SM, el número de \textit{warps} activos y elegibles, la utilización de los Tensor Cores, el tráfico a través de memoria global y caché L2, y las principales causas de espera de los \textit{warps}.

En particular, estas pruebas permiten diferenciar situaciones en las que el rendimiento está limitado por el ancho de banda de memoria de aquellas en las que el principal problema es la falta de paralelismo disponible para ocultar la latencia. También permiten analizar el efecto sobre el rendimiento de parámetros de implementación como el número de \textit{warps} por bloque, el número de buffers utilizados en memoria compartida o el consumo de registros.

Los tests de \textit{profiling} se mantienen separados de los tests temporales. Esto es necesario debido a que la instrumentación realizada por Nsight Compute puede requerir múltiples ejecuciones o \textit{replays} del kernel y, por tanto, los tiempos registrados durante una sesión de \textit{profiling} no se utilizan como medidas directas de rendimiento.

\section{Organización de las pruebas}

La separación entre validación funcional y pruebas de rendimiento permite utilizar cada conjunto con un objetivo diferente durante el desarrollo. El test funcional se ejecuta después de realizar modificaciones en los kernels para comprobar que se mantiene la correctitud, mientras que los tests temporales y de \textit{profiling} se utilizan posteriormente para determinar si dichas modificaciones producen una mejora real de rendimiento.

Esta metodología permite que una optimización sea considerada válida únicamente cuando mantiene la correctitud numérica y produce una mejora medible bajo las mismas condiciones experimentales.

